% Generated by scripts/polish_appendix_tables.py; edit the formatter, not this file.
\begin{table}[H]
\caption{\textbf{Within-family difficulty settings.} L1--L8 specify each family\textquotesingle s computation and numerical representation.}
\label{tab:difficulty-definitions}
\centering\tablefont
\setlength{\tabcolsep}{3pt}
\begin{tabular}{@{}p{.10\linewidth}p{.40\linewidth}p{.43\linewidth}@{}}
\toprule
\textbf{Family} & \textbf{L1--L4} & \textbf{L5--L8} \\
\midrule
Comparison & L1: one-digit positive integers; L2: two-digit positive integers; L3: two-digit negative integers; L4: signed integers. & L5: hundredths; L6: rational fractions; L7: scientific notation; L8: percentage versus decimal. \\
\addlinespace[3pt]
Sum & L1: 1+1 digit; L2: 2+1 digits; L3: 3+2 digits; L4: 4+3 digits. & L5: two 3-digit scaled operands (one decimal place); L6: three 2-digit integers; L7: 5+4 digits; L8: three scaled operands (3,3,1 digits; one decimal place). \\
\addlinespace[3pt]
Difference & L1: 1$-$1 digit; L2: 2$-$1 digits; L3: 3$-$2 digits; L4: 4$-$3 digits. & L5: 3$-$2 scaled digits (one decimal place); L6: 6$-$4 digits; L7: $A-B-C$ (4,3,2 digits); L8: $A-(B-C)$ (4,3,2 digits). \\
\addlinespace[3pt]
Expression & L1: addition; L2: subtraction; L3: multiplication; L4: $A+B\times C$. & L5: $(A+B)\times C$; L6: $A\times B+C\times D$; L7: $(A-B)\times(C+D)$; L8: $A+B\times(C-D)$. The long factor appends one term. \\
\addlinespace[3pt]
Product & L1: $1\times1$ digit; L2: $2\times1$ digits; L3: $2\times2$ digits; L4: $3\times1$ digits. & L5: $2\times1$ scaled digits (one decimal in A); L6: $2\times2$ scaled digits (one decimal each); L7: $A\times B\times C$; L8: $A\times(B+C)$, with one decimal in A. \\
\addlinespace[3pt]
Percentage & L1: $10\%$ of N; L2: $25/50\%$; L3: $5/20\%$; L4: $15/30\%$. & L5: fractional percentage and decimal N; L6: increase N by $P\%$; L7: recover the original before an increase; L8: successive increase and decrease. \\
\bottomrule
\end{tabular}
\appnote{Digit counts describe unscaled integer operands; decimal settings are stated separately. Each family--level cell has 30 cases, split 15/15 between the factor values in Table~\ref{tab:families}.}
\end{table}
